% Overleaf: select XeLaTeX. UTF-8. Self-contained: no external figures or .bib.
% This is a research log, not a final paper or a claim of completed circuit discovery.
\documentclass[11pt,a4paper]{article}
\usepackage[margin=2.3cm]{geometry}
\usepackage{fontspec}
\usepackage{amsmath,amssymb,booktabs,longtable,array}
\usepackage{xcolor}
\usepackage{tikz}
\usetikzlibrary{positioning,arrows.meta}
\usepackage[unicode,colorlinks=true,linkcolor=blue!50!black,urlcolor=blue!60!black]{hyperref}
\usepackage{polyglossia}
\setdefaultlanguage{hebrew}
\setotherlanguage{english}
\setmainfont{FreeSerif}
\newfontfamily\hebrewfont[Script=Hebrew]{FreeSerif}
\newfontfamily\englishfont{TeX Gyre Termes}
\newfontfamily\hebrewfonttt{DejaVu Sans Mono}
\setmonofont{DejaVu Sans Mono}
\newcommand{\en}[1]{\textenglish{#1}}
\newcommand{\file}[1]{\textenglish{\nolinkurl{#1}}}
\setlength{\parindent}{0pt}
\setlength{\parskip}{5pt}
\setlength{\emergencystretch}{3em}
\renewcommand{\arraystretch}{1.2}
\title{מראשי אינדוקציה סמנטיים למעגלים של הסקה מרובת צעדים\\\large יומן מחקר, החלטות, ניסויים ומגבלות}
\author{דניאלה סימונובסקי ומקס גולובקוב}
\date{תיעוד ביניים לפי השיחה וקובצי הפרויקט הזמינים}

\begin{document}
\maketitle
\begin{abstract}
מסמך זה מתעד את התפתחות הפרויקט מהצעת המחקר ועד לפיילוטים ההתנהגותיים האחרונים שנמצאו בתיקיית הפרויקט. נקודת המוצא הייתה השאלה האם עיבוד קשרים סמנטיים מרובי צעדים נשען על ראשי קשב בודדים או על מעגל מבוזר. בהמשך גיבשנו מסגרת ניתוח מכניסטית, בחרנו משפחת מודלים עם גישה למשקולות ולנקודות ביניים באימון, בנינו תשתית ניסויים ודאטה סינתטי, והרצנו סדרת מבחני יכולת. הפיילוטים חשפו הטיות תשובה, תלות בפורמט ובהדגמות, עמימות לוגית וקיצורי דרך בדאטה. אלה הובילו לשינויים מתועדים במשימה ובבקרות. נכון לתיעוד זה, קיימות תוצאות התנהגותיות ותשתית הערכה; טרם זוהה או אומת מעגל סמנטי באמצעות התערבויות פנימיות.
\end{abstract}
\tableofcontents
\newpage

\section{מטרת התיעוד ומעמד הראיות}
המסמך נועד לשמש זיכרון מחקרי לכתיבת העבודה הסופית: מה רצינו לבדוק, מה בנינו, מה הרצנו, מדוע שינינו כיוון, ומה עדיין אינו מבוסס. הוא אינו מאמר תוצאות סופי. חלק מהניסוחים בהצעה ובמסמכי התכנון היו שאפתניים יותר מן הראיות שנאספו; כאן מובחנים במפורש תכנון, מימוש ותוצאה.

המקורות לשחזור הם הצעת המחקר, השיחה, מסמכי התכנון, מחוללי הדאטה, מניפסטים, דוחות ורשומות תוצאות. תוצאות גרסאות ההשלמה המאוחרות נמצאו בתיקייה גם מעבר לנקודת העצירה של השיחה המצוטטת. הן נכללות כדי לתעד את המצב בפועל. אין להסיק תאריך ביצוע מתוך מספר \en{seed} או הערת תאריך בקוד; רצף השלבים להלן הוא לוגי וגרסאי, ולא יומן תאריכים מאומת.

שלוש רמות הראיות במסמך:
\begin{enumerate}
\item \textbf{בוצע ומתועד:} קוד, דאטה או תוצאות שנשמרו בפועל; מספרים שחושבו מן הרשומות או מדוחות קודמים מזוהים.
\item \textbf{תוכנן או נדון:} שיטת ניתוח או ניסוי שהוגדרו אך אין עבורם תוצאה פנימית של המודל.
\item \textbf{השערה או מגבלה:} הסבר אפשרי לתוצאה שאינו ממצא סיבתי מוכח.
\end{enumerate}

\section{נקודת ההתחלה: הצעת המחקר}
כותרת ההצעה הייתה \en{From Heads to Circuits -- Identifying Semantic Induction Circuits (SICs) for multi-hop Reasoning in LLMs}. שאלת המחקר המקורית הייתה: כיצד מודלי שפה מעבדים מבחינה מכניסטית קשרים סמנטיים מרובי צעדים? האם ראש \en{SIH} בודד מספיק, או שהפעולה דורשת מעגל של רכיבים לאורך כמה שכבות?

המקרה המנחה היה הסקה טרנזיטיבית:
\[
 A\subseteq B,\qquad B\subseteq C\qquad\Longrightarrow\qquad A\subseteq C.
\]
כאן הסימון מתאר הכלת קבוצות, כפי שבמשפטי \en{All A are B}; הוא אינו דורש הנחה של הכלה ממש. שליפת $A\subseteq B$ היא בקרת צעד אחד, ולא הפעלת טרנזיטיביות כשלעצמה.

ההצעה כללה שימוש בסדרות מודלים כגון \en{Pythia}, \en{OLMo} או \en{LMEnt}, ניתוח לאורך נקודות אימון, \en{activation patching}, \en{path patching}, הרחבת \en{Relation Index} למדד מעגלי והערכת \en{faithfulness} של תת־גרף דליל. ב־\en{bAbI} וב־\en{CLUTRR} ראינו מקורות אפשריים לתבניות ולמשימות; לא תיעדנו הרצה של המאגרים המקוריים כפיילוט שבוצע.

הנחות ההיתכנות הראשוניות היו מודל קטן, גישת \en{GPU} באוניברסיטה ודאטה סינתטי עם ישויות מומצאות. בהמשך התברר שהיכולת ההתנהגותית עצמה אינה מובטחת, ולכן היא נעשתה שער כניסה למחקר המכניסטי. גם ההנחה ששמות מומצאים מבטיחים הסתמכות טהורה על ההקשר רוככה: הם מפחיתים אפשרות לשלוף עובדות מוכרות על הישויות, אך אינם מבטלים ידע קודם על שפה, יחסים, מבנה המשימה או רמזים טוקניים.

\section{בניית הבסיס התאורטי}
\subsection{שני מאמרי אנתרופיק: קריאה מעמיקה}
לפי תיעוד העבודה, קראנו לעומק את \en{In-context Learning and Induction Heads} ואת \en{A Mathematical Framework for Transformer Circuits} \cite{olsson,framework}. למדנו להפריד בין דפוס קשב לבין המידע שהראש מעביר, ולחשוב על הזרם השאריתי כאמצעי שבו רכיבים מתקשרים.

במנגנון האינדוקציה הקלאסי התבנית היא בקירוב
\[
\ldots,A,B,\ldots,A \quad\longrightarrow\quad B.
\]
במימוש הפשוט בשתי שכבות, ראש מוקדם מעביר מידע מן הטוקן הקודם, וראש מאוחר משתמש בו כדי לאתר המשך מתאים ולהעתיקו. השם \en{induction head} מתייחס בדרך כלל לראש המאוחר; המנגנון המלא כולל יותר מראש אחד. לכן אין לזהות אוטומטית ``ראש'' עם ``מעגל''.

המסגרת המתמטית סיפקה הבחנה בין \en{QK}, הקובע לאן מופנה הקשב, לבין \en{OV}, הקובע מה מועבר משם. בנוסף עסקנו בהרכבת ראשים דרך \en{Q-, K-, V-composition}: פלט של רכיב מוקדם יכול להשפיע על שאילתות, מפתחות או ערכים של רכיב מאוחר. זה סיפק שפה מדויקת יותר לשאלה כיצד שתי פעולות מקומיות עשויות להתחבר לפתרון של שתי קפיצות.

עלה גם כיוון של תלות בהקשר ארוך מטוקן קודם יחיד: הבחנה בין העתקת תבניות ארוכות יותר לבין הרכבת קשר סמנטי. הוא נשאר שאלת רקע, ולא תוצאת ניסוי שלנו.

\subsection{המאמר המרכזי על SIH}
התבססנו ישירות על Ren ואחרים \cite{ren}. המאמר בוחן ראשים שמקשרים בין ישויות ביחסי ידע או בתלויות תחביריות: כאשר ראש מתייחס לטוקן של ישות המקור, הקריאה של פעולת ה־\en{OV} שלו עשויה לקדם טוקן הקשור אליה. המדד \en{RI} משמש לאפיון דפוסים אלה.

הבחנה חשובה שעלתה בדיון היא ש־\en{SIH} אינו בהכרח ראש שעבר קודם סינון לפי מבחן האינדוקציה הקלאסי. לכן תכננו לבדוק בנפרד ציוני העתקה והתאמת קידומת, ולא להניח שהכינוי ``סמנטי'' מוכיח מימוש של מעגל האינדוקציה המקורי.

גם ההשוואה ההתנהגותית למאמר דורשת זהירות. שם נבחנים בין השאר יחסי ידע ותלויות תחביריות בטקסט, והניסויים אינם זהים לפתרון שרשרת חדשה בין שמות מומצאים. מדד \en{RI}, דיוק במשימות \en{ICL}, ודיוק במבחן הסינתטי שלנו אינם אותו משתנה. לפיכך ביצועי מודל קטן יותר במאמר אינם קו בסיס ישיר לביצועי \en{Pythia} אצלנו.

\subsection{ספרות נוספת וכיצד היא נכנסה לתכנון}
עברנו על כיוונים ומקורות נוספים, ונאספה ספריית מאמרים. אין תיעוד שמאפשר לייחס אותה רמת קריאה לכל קובץ; שני מאמרי אנתרופיק והמאמר המרכזי הם בסיס הקריאה העמוקה, והאחרים הם מקורות מתודולוגיים או ספרות עזר.

\begin{itemize}
\item \en{Interpretability in the Wild}: דוגמה למעבר מזיהוי רכיבים למיפוי מעגל ולהערכת תרומתו \cite{ioi}.
\item \en{How to Use and Interpret Activation Patching} וכן \en{Towards Best Practices of Activation Patching in Language Models}: בחירת התערבות, מדד תוצאה וקלטי מקור/יעד \cite{patchguide}.
\item \en{Towards Automated Circuit Discovery}, גילוי מעגלים מודע מיקום ועבודות על \en{faithfulness}: דרכים להרחיב חיפוש ולהעריך תת־גרף מעבר לחפיפה בין רשימות ראשים.
\item \en{Transformer Feed-Forward Layers Are Key-Value Memories} ו־\en{Toy Models of Superposition}: תזכורת לכך שמידע רלוונטי יכול להיות ב־\en{MLP}, במאפיינים מבוזרים או בסופרפוזיציה, ולא רק בראש קשב בודד.
\item מאמרים על התפתחות ראשים סטטיסטיים, דינמיקת למידה ו־\en{grokking}: רקע לכיוון העתידי של מעקב אחר נקודות אימון.
\end{itemize}
אין להסיק מהימצאות מאמר בתיקייה שבוצע שחזור שלו. בפרט, לא נמצא תיעוד לניסוי שלנו שמאמת מאמר המשך מסוים ל־\en{SIH}.

\section{המתודולוגיה שגיבשנו: כמה מקורות ראיה משלימים}
רצינו לא להסתמך על מדד יחיד. ``שיטות בלתי תלויות'' הובן כאן כראיות מסוגים שונים; לא נטענה עצמאות סטטיסטית ביניהן, שכן הן עשויות להשתמש באותם קלטים ובאותם רכיבים.

\subsection{שער התנהגותי לפני פרשנות פנימית}
הסדר שנבחר הוא: להוכיח תחילה שיש התנהגות יציבה להסביר, ורק אחר כך לחפש את המנגנון. תשובה נכונה על דוגמאות נבחרות אינה מספיקה: דרושה גם רגישות לשינוי העובדות ועמידות לשינויים לא רלוונטיים. אם קשר ישיר נכשל, אין טעם לפרש ראש כמרכיב של הסקה מורכבת על סמך אותו מבחן בלבד.

\subsection{RI וקריאות מקומיות: סינון מועמדים}
תכננו להתחיל משחזור איכותי של \en{RI} על יחסים פשוטים, עם בקרות של יעד אקראי וראש אקראי. זהו מדד תצפיתי לתאימות בין פעולת ראש לבין קשר מוגדר; הוא אינו שקול לטענה שהראש ``מבין'' את הקשר, ואינו מוכיח נחיצות או מספיקות לפלט.

עלתה הצעה להרחיב את המדד באמצעות התניה: למשל לבדוק אם פעילות של $H_2$ הקשורה לקשר השני מתרחשת יותר כאשר $H_1$ מקדם את הקשר הראשון. זו הצעה שנדונה, לא מדד סופי שמומש ואומת. מכפלת ציוני שני ראשים או מתאם מותנה לבדם לא מוכיחים מעבר מידע ביניהם; שכבה, זהות טוקנים וקושי הדוגמה יכולים להסביר את שניהם.

\subsection{התערבויות ברכיבים ובמסלולים}
\en{Activation patching} הוא ניסוי התערבותי בתוך חישוב המודל: מחליפים הפעלה בערך מקלט אחר ובודקים שינוי בתוצאה. לכן אין לתאר אותו כקורלציה בלבד. עם זאת, המסקנה הסיבתית מוגבלת להתערבות, לקלטים, למיקום ולמדד שנבחרו; ערך מלאכותי או ערבוב מצבים לא טבעי עלולים להקשות על הפרשנות.

לצורך ציון רציף אפשר להגדיר
\[
 m(x)=\log P(y_{\mathrm{target}}\mid x)-\log P(y_{\mathrm{alternative}}\mid x).
\]
בהינתן קלט מקור תקין $x_c$, קלט ששונה $x_b$ והתערבות ברכיב $z$, נבחן את השינוי ב־$m$. אם מציגים שיעור שיקום מנורמל,
\[
 R_z=\frac{m(x_b;z\leftarrow z(x_c))-m(x_b)}{m(x_c)-m(x_b)},
\]
יש להגדיר את אותו יעד בשני הקלטים, לדווח גם את השינוי הלא מנורמל ולהיזהר כשהמכנה קטן. אלה עקרונות תכנון; טרם חושבו בפרויקט תוצאות \en{patching} כאלה.

למיפוי מעגל, לא מספיק למצוא שני ראשים שכל אחד מהם משנה את התשובה. תכננו לבדוק מסלול מכוון מרכיב מוקדם לרכיב מאוחר באמצעות \en{path patching}, להפריד תרומה דרך \en{Q/K} מתרומה דרך \en{V}, ולבחון האם השפעת הרכיב המוקדם עוברת דרך הרכיב המאוחר. שינוי במשקולות הקשב אחרי החלפת הפעלה אינו שינוי במשקולות המודל עצמן.

\subsection{בדיקות מעגל שלם והשערות חלופיות}
תכננו לבחון נחיצות באמצעות הסרה, מספיקות יחסית באמצעות שימור תת־גרף, ו־\en{faithfulness} ביחס למודל המלא. צריך לקבוע את בסיס ההסרה ואת סף השיקום מראש; אין מספר קסם שמוכיח מעגל. יש לבדוק קלטים חדשים, רכיבים חלופיים, \en{MLP} ומסלולים מקבילים, ולא לכפות הסבר של שני ראשים מפני שהוא תואם את מספר הקפיצות.

בדיקות אינדוקציה קלאסית, בקרות מידע מקומי ובדיקות לאורך נקודות אימון נשארו מסלולים משלימים מתוכננים. לא תועד יישום של \en{causal scrubbing}, חיפוש מעגל אוטומטי או סריקת נקודות אימון.

\begin{figure}[ht]
\centering
\begin{english}
\begin{tikzpicture}[node distance=7mm,box/.style={draw,rounded corners,align=center,text width=10.5cm,minimum height=9mm},>=Stealth]
\node[box,fill=blue!8] (a) {Behavioral validation: direct, two-hop, counterfactuals, robustness};
\node[box,below=of a] (b) {Candidate components: RI, attention, local output diagnostics};
\node[box,below=of b] (c) {Node interventions: activation patching and ablation};
\node[box,below=of c] (d) {Directed paths: mediation through Q/K/V and path patching};
\node[box,below=of d] (e) {Circuit validation: faithfulness, controls, held-out families};
\draw[->] (a)--(b); \draw[->] (b)--(c); \draw[->] (c)--(d); \draw[->] (d)--(e);
\end{tikzpicture}
\end{english}
\caption{סדר העבודה המתוכנן. העבודה שבוצעה עד כה מתמקדת בשלב הראשון ובתשתית הנדרשת להמשך.}
\end{figure}

\section{בחירת המודל והמשאבים}
בחרנו להתחיל ב־\en{Pythia-1B}: מודל פתוח קטן יחסית, עם אפשרות לגשת למשקולות ולנקודות אימון, ובהמשך לעבור באותה משפחה לגודל גדול יותר. הבחירה ב־\en{Pythia-6.9B} התקבלה במפורש לאחר הפיילוט הקטן; לא בחרנו בפועל ב־\en{InternLM2-7B}. עצם הזכרת \en{OLMo} וחלופות אחרות אינה שקולה להרצתן.

שני המודלים הופעלו כמודלי בסיס באנגלית. הדגמות \en{ICL} נוספו להקשר, בלי עדכון משקולות. לא בוצעו אימון מאפס או \en{fine-tuning}. נקודות המודל שננעלו מקומית הן:
\begin{english}\small
\begin{itemize}
\item Pythia-1B: \texttt{f73d7dcc545c8bd326d8559c8ef84ffe92fea6b2}.
\item Pythia-6.9B: \texttt{c0e3eee36dc47af0c49f361c74cfe459c09f7f23}.
\end{itemize}
\end{english}

המשתמש דיווח על גישה ל־\en{GPU} ולאחסון ב־\en{TAU Slurm}, ועל חלון עבודה של כשלושה שבועות בעת התכנון. תכננו תשתית שניתן להעביר גם ל־\en{Colab} או \en{RunPod}, אך התוצאות המתועדות כאן הן מן ההרצות שהוחזרו מהאשכול. אין תיעוד לניסוי מקביל שבוצע אצל ספק אחר.

הפרדנו בין הורדת קובצי המודל לבין הרצה וטעינתם לזיכרון. בדיקת התנהגות מקומית על השרת דורשת גישה למשקולות, לא רק פתיחת כרטיס המודל בדפדפן. העברת החבילות נעשתה בין המחשב האישי לשרת; אין להסיק ממסמך התקנה מוקדם שכל המשקולות הורדו תמיד באותו מסלול.

\section{תשתית הדאטה וההרצה}
\subsection{גרף לפני טקסט}
בגרסאות הראשונות בנינו את העובדות כגרף ורק אחר כך הפכנו אותן למשפטים. כך ניתן לחשב תשובת אמת ומסלול הוכחה, לבדוק שאין מסלול קצר יותר, ולבנות שינוי מבוקר. שמרנו מזהה משפחה ושאלה, עובדות, שאילתה, תשובה, מספר צעדים, מסיחים, סדר, \en{seed}, טווחי ישויות וקישור לגרסת בסיס. בזמן ההרצה נשמרו גם מזהי טוקנים, פלט וציונים.

אין להניח ששם הוא טוקן יחיד. לכן נבנתה תמיכה בטווחי תווים וטוקנים ובניקוד תשובות רב־טוקניות. בגרסאות מאוחרות הסכמה השתנתה בחלקה; יש לשמר את הפרומפט המדויק ולא להסתפק בשם הגרסה.

\subsection{ארבעת מבחני היכולת}
\begin{enumerate}
\item \textbf{קשר ישיר:} שליפה או זיהוי של קשר המופיע בעובדה אחת.
\item \textbf{שני צעדים:} שילוב שתי עובדות דרך ישות ביניים.
\item \textbf{שינוי הקשר:} ניתוק שרשרת או ניתוב מחדש, בהתאם לגרסת המשימה.
\item \textbf{עמידות:} שינוי סדר, הוספת מסיחים והחלפת שמות.
\end{enumerate}
משמעות המבחן השלישי השתנתה במהלך העבודה. בכן/לא ניתוק דרש \en{No}; בהשלמת שמות השינוי יצר יעד תקף אחר. אין לאחד את הדיוקים כאילו מדובר באותו מבחן.

\subsection{מעטפת שחזור ובקרת ריצות}
נכתבו מחוללי דאטה, סקריפטי הרצה וניתוח, הגשה ל־\en{Slurm}, בדיקות עשן, מניפסטים וטביעות \en{SHA-256}. נבנתה חלוקה למספר עותקי מודל לפי מספר ה־\en{GPU} שהוקצה, כאשר משפחה נשארת כיחידה אחת. עבור המודל הגדול תועד גם פיצול מודל על שני כרטיסי \en{RTX 2080 Ti} בני \en{11GB}; אין להניח שזה החומרה של כל עבודה מאוחרת.

מיזוג תוצאות בודק כיסוי מזהים והתאמה לדאטה. חידוש ריצה מותנה באותם קבצים והגדרות; שינוי מודל, דאטה או קוד מחייב שם ריצה חדש. נשמרו תוצאות גולמיות לצד דוחות, כדי לבדוק קלט ופלט ולא להסתפק בטבלת דיוק.

\subsection{תקלות מספריות וסביבתיות}
לפני הריצות המלאות היו כשלי התקנה, טעינה וחישוב, ותועדו לוגים לאבחון. במודל הגדול הופיעו ערכים לא סופיים בנתיב קשב שנבדק. מעבר ל־\en{SDPA MATH}, עם משקולות \en{FP16} ובדיקות סופיות של ציוני התשובות, אפשר את הריצות התקינות. אין בכך זיהוי ודאי של הפעולה המספרית הראשונה שיצרה את התקלה.

נדון גם אימות מדגם ב־\en{FP32} וטיפול בתיקו: אם שני ציונים מתעגלים לאותו ערך, הכרעה לפי סדר האפשרויות יכולה ליצור הטיה. זו הצדקה לבדיקה, לא תיעוד לכך שבוצעה הרצה מלאה נוספת ב־\en{FP32}. הפרדנו בין תקלות טכניות לבין חולשה התנהגותית; הצלחת בדיקת עשן אינה הוכחה ליכולת הסקה.

\section{מדדי ההערכה והמשמעות שלהם}
בהשלמת שמות, \en{free accuracy} הוא שיעור ההפקות החופשיות שבהן השם שחולץ שווה ליעד. \en{Candidate accuracy} משווה בין שתי השלמות מוגדרות מראש. עבור שם $y$ עם טוקנים $t_1,\ldots,t_k$ נצבר הציון
\[
 s(y\mid x)=\sum_{i=1}^{k}\log P(t_i\mid x,t_{<i}).
\]
בקוד נכללים הרווח המוביל והנקודה לפי הטוקניזציה בפועל. נשמר גם ממוצע לטוקן כאבחון, בלי להחליף בדיעבד מדד משום שהוא נותן תוצאה נוחה יותר.

הפקה חמדנית והעדפת הרצף השלם אינן אותה פעולה, ולכן שני המדדים יכולים לחלוק. קו בסיס של $50\%$ שייך לבחירה בין שתי אפשרויות מאוזנות; הוא אינו קו בסיס כללי להפקה מתוך אוצר המילים. בגרסאות השלמה מוקדמות שתי האפשרויות מופיעות גם בפרומפט החופשי; בגרסאות השפה הטבעית המאוחרות הן קיימות רק לצורך הניקוד.

נוסף לדיוק, בדקנו שיעורי תוויות, התאמה לפורמט, תיקו, שמות ביניים, מיקום אפשרויות, נגישות בגרף והצלחה בשני בני זוג. אי־ודאות צריכה להימדד לפי משפחות: $1{,}200$ גרסאות תלויות אינן $1{,}200$ תצפיות עצמאיות.

\section{רצף הפיילוטים: מה שונה ומדוע}
\subsection{שלב א: Pythia-1B, שאלות כן/לא}
הדאטה כלל שמות כגון \en{group000a}, ארבע הדגמות קבועות ומבחני כיוון וקשר. \en{No} פירושו ``לא נובע מהעובדות'', ולא ``ההפך הוכח''. הוכן גם תנאי ללא הדגמות; תוצאת ההרצה המלאה המזוהה כאן היא ארבע הדגמות.

ב־$1{,}200$ התוצאות שנשמרו, הדיוק החופשי בקשר הישיר היה $53.5\%$, ובשני צעדים $61.0\%$. הפירוט לפי תווית חשף העדפת \en{Yes}: בקשר הישיר $88\%$ הצלחה בחיוביות לעומת $19\%$ בשליליות. בשבירת החוליה הראשונה והשנייה התקבלו בהתאמה $20\%$ ו־$17\%$ דיוק חופשי. התוצאות לא הצדיקו מעבר ישיר לפירוש מעגל.

החלטנו לבחון מודל גדול יותר, לשפר את מובחנות השמות ולהגדיל את מספר ההדגמות. לא קבענו שההבדל בין השמות הוא שגרם לכישלון; זה היה גורם אפשרי שאפשר לשנות.

\subsection{שלב ב: Pythia-6.9B, כן/לא עם 4 ו־12 הדגמות}
הוחלפו השמות בשמות מובחנים כגון \en{wugs} ו־\en{daxes}. לאחר פתרון תקלות ההרצה התקבלו שתי ריצות מלאות, $1{,}200$ קלטים בכל אחת. בארבע הדגמות המודל הפיק \en{No} בכל השאלות. בשתים־עשרה התקבלו $1{,}154$ הפקות \en{No}; הדיוק החופשי הכולל היה $698/1200$ לעומת $700/1200$ בארבע הדגמות.

בדאטה היו $700$ שליליות ו־$500$ חיוביות. לכן תמיד־\en{No} משיג $58.33\%$ כולל, ואף $100\%$ בשבירת שרשרת שכל שאלותיה שליליות. זו הייתה דוגמה ברורה לכך שדיוק כולל מסתיר כשל בסיסי.

\subsection{שלב ג: תיקון ההדגמות בכן/לא}
בגרסת \en{prompt-v3} של כן/לא גיוונו הדגמות בין משפחות, איזנו תוויות ותווית אחרונה, והתאמנו טוב יותר את המבנה לדאטה. יש להבחין בין שם גרסה זה לבין \en{completion-v3} המאוחר, שהוא משימה אחרת.

המודל נשאר כמעט תמיד עם \en{No}: בארבע הדגמות כל $1{,}200$ התשובות היו \en{No}; בשתים־עשרה היו $1{,}194$ כאלה ושש \en{Yes}. הדיוק הכולל היה $700/1200$ ו־$704/1200$. בקשר ישיר ובשני צעדים נשאר $50\%$. מדדי דירוג רציפים לא הצביעו על אות חזק של הסקה שמוסתר רק על ידי סף התשובה. חיקוי של התווית האחרונה לבדו לא הסביר את הדפוס.

\subsection{שלב ד: completion-v1, מעבר להשלמת שם}
לבקשת המשתמש עברנו מכן/לא להשלמת ישות, למשל \en{Therefore, all korvas are ...}. נבנו $1{,}200$ שאלות עם $0$, $4$ ו־$12$ הדגמות. בשלב זה הוצגו שני מועמדים, והשם שבאמצע השרשרת הוצא מן האפשרויות.

\begin{center}\small
\begin{tabular}{lrrr}
\toprule
מבחן & 0 הדגמות & 4 הדגמות & 12 הדגמות\\
\midrule
קשר ישיר & \en{56.0 / 70.5} & \en{79.5 / 91.0} & \en{77.0 / 88.5}\\
שני צעדים & \en{17.5 / 50.0} & \en{28.0 / 71.5} & \en{35.0 / 75.0}\\
ניתוב מחדש & \en{22.5 / 58.0} & \en{20.0 / 59.0} & \en{26.5 / 55.0}\\
עמידות & \en{18.2 / 55.5} & \en{20.8 / 67.0} & \en{29.7 / 65.0}\\
\bottomrule
\end{tabular}
\end{center}
בכל תא כאן ובהמשך: \en{free / candidate}, באחוזים. במבחנים הראשונים $200$ קלטים לכל מבחן ובאחרון $600$.

הפורמט אפשר הצלחות רבות יותר, אך האבחון חשף שתי בעיות:
\begin{enumerate}
\item \textbf{שם הביניים:} ב־12 הדגמות ובשני צעדים, $70/200$ הפקות היו היעד הנכון, $100/200$ היו שם הביניים ו־$30/200$ המסיח. בשרשרת $A\to B\to C$, גם השלמת $B$ אחרי ``כל $A$ הם'' נכונה לוגית, אף שאינה האפשרות המבוקשת.
\item \textbf{רמז שכיחות:} וידאנו ששני המועמדים מופיעים בעובדות, אך לא איזנו את מספר הופעותיהם. כלל שבוחר את השם השכיח יותר וניחוש אחיד בתיקו השיג $83.3\%$ בתוחלת בכלל הדאטה; במבחן שני הצעדים הבסיסי הוא השיג $100\%$.
\end{enumerate}
זו מגבלה של הדאטה שבנינו, לא הוכחה שהמודל משתמש בכלל השכיחות. גם נמצאה תלות במיקום האפשרויות: ללא הדגמות הועדפה הראשונה, ועם הדגמות השנייה. אין לייחס את השיפור להסקה בלי לנטרל רמזים אלה.

\subsection{דוגמה מתועדת: הצלחה ושינוי התשובה}
\begin{english}\small
\begin{verbatim}
All tivaks are snorps.
All vromps are ruspins.
All jastles are kelbrins.
All ruspins are kelbrins.
Possible completions: snorps or kelbrins.
Therefore, all vromps are
\end{verbatim}
\end{english}
בשאלה \file{003/twohop/base/0} התשובה הנכונה וההפקה היו \en{kelbrins}. בגרסת \file{003/broken/first/0} הוחלף המשפט השני ב־\en{All vromps are tivaks}. התשובה הנכונה נעשתה \en{snorps}, אך המודל המשיך לבחור \en{kelbrins}. השם הזה עדיין הופיע פעמיים והיה האפשרות השנייה. לכן הדוגמה מתיישבת עם כמה הסברים חלופיים ואינה מבדילה ביניהם.

בגרסה זו, ב־12 הדגמות, הצלחה חופשית גם במקור וגם לאחר שינוי הקשת הראשונה הייתה $3\%$, ולאחר שינוי השנייה $13\%$. מדדי \en{candidate} המקבילים היו $19\%$ ו־$45\%$. אין לבלבל מדד הצלחה בשני קלטים עם דיוק בקלט ששונה לבדו.

\subsection{שלב ה: completion-v2, עומק מפורש ובקרות מאוזנות}
באישור המשתמש הוגדרה משימה מפורשת של יעד במסלול הקצר ביותר באורך צעד אחד או שניים. ההוראה הופיעה גם בהדגמות, ושם הביניים כבר לא היה יעד תקף לשאלת שני צעדים גם בלי רשימת האפשרויות.

בנוסף איזנו שכיחות ותפקיד תחבירי של המועמדים, הוספנו זוגות עם סדר אפשרויות הפוך, ושינינו את בקרת הניתוב מחדש להחלפת יעדים בין שתי עובדות כדי לשמר ספירות. גודל כל תנאי נשאר $1{,}200$: $50$ משפחות כפול $12$ שאלות כפול שני סדרי אפשרויות. לכן היו $600$ שאלות תוכן, ולא $1{,}200$ שאלות עצמאיות. בדיקות הבנייה עברו; כלל השכיחות ירד ל־$50\%$ בתוחלת, וכללי מיקום בעובדות השיגו בערך $48$--$53\%$ במבחנים הראשיים.

נמצאו בתיקייה תוצאות מלאות לארבע ולשתים־עשרה הדגמות:
\begin{center}\small
\begin{tabular}{lrr}
\toprule
מבחן & 4 הדגמות & 12 הדגמות\\\midrule
קשר ישיר & \en{50.5 / 50.5} & \en{56.0 / 56.0}\\
שני צעדים & \en{49.0 / 49.5} & \en{52.0 / 52.0}\\
ניתוב מחדש & \en{51.5 / 53.0} & \en{46.5 / 46.5}\\
עמידות & \en{47.7 / 48.0} & \en{50.2 / 50.2}\\\bottomrule
\end{tabular}
\end{center}
הגרסה לא פתרה את בעיית היכולת. התוצאה תואמת רגישות לפורמט או קושי בדרישה המטא־לשונית, אך מאחר ששונו גם המבנה והבקרות, לא הוכח שהניסוח לבדו גרם לירידה. זו תוצאה שלילית מתועדת, לא סיבה למחוק את הגרסה מן ההיסטוריה. הוכן גם דאטה ללא הדגמות; לא נמצאה עבורו ריצה מלאה בגרסה זו בבדיקה הנוכחית.

\subsection{שלב ו: completion-v3, יחסים בשפה טבעית}
בקבצים המאוחרים עברנו לניסוחים טבעיים בלי שורת אפשרויות ובלי הוראת ``מסלול קצר ביותר''. עוגן המשימה הוא יחס מורכב בעל שם: אם וסבתא. אם $B$ היא אמו של $A$ ו־$C$ היא אמו של $B$, השאלה על סבתו של $A$ מצביעה על $C$; $B$ אינה תשובה שקולה לשאלה. זה יחס הרכבה, לא טרנזיטיביות של ``אם'' עצמה.

נוספו גם עולמות הכלה ומיקום. שם העמימות של ישות ביניים לא נעלמת לחלוטין בהפקה חופשית, ולכן דרוש סיווג נפרד שלה. $68$ מתוך $100$ משפחות הן יחסי משפחה, $16$ הכלה ו־$16$ מיקום: הממוצע הכולל אינו ממוצע שווה משקל של שלושת העולמות.

בכל משפחה שש שאלות ושתי גרסאות סדר עובדות, יחד $1{,}200$ קלטים. נשמרו מועמדים לניקוד אך הם אינם מוצגים בפרומפט. נבדקו איזון אזכורים ואורכי טוקנים של המועמדים. הוכנו גם $40$ בקרות נפרדות של עובדה מפורשת וחזרה מילולית על הדגמה. לא נמצאה בבדיקה הנוכחית תוצאת הרצה של בקרות אלה, ולכן הן מסומנות כהכנה בלבד. כישלון בהן יצריך אבחון; אין להכריז אוטומטית שכל סטייה מ־$100\%$ היא באג.

בגרסת \en{v3} הראשונית התקבלו:
\begin{center}\small
\begin{tabular}{lrr}
\toprule
מבחן & 4 הדגמות & 12 הדגמות\\\midrule
קשר ישיר & \en{84.0 / 86.5} & \en{85.5 / 87.0}\\
שני צעדים & \en{61.0 / 71.0} & \en{59.5 / 67.0}\\
ניתוב מחדש & \en{40.0 / 46.0} & \en{46.5 / 50.5}\\
עמידות & \en{42.7 / 60.2} & \en{45.3 / 63.5}\\\bottomrule
\end{tabular}
\end{center}

\subsection{שלב ז: completion-v3.1, איזון עומק וסדר בהדגמות}
השוואת הקוד המקומי לקוד בחבילת ההעלאה הראתה שינוי נוסף. ב־\en{v3} מספר הקפיצות וסדר משפחות העובדות בהדגמות היו מצומדים. ב־\en{v3.1} נוצרו ארבעת הצירופים של עומק וסדר, ונבחר מספר שווה מכל צירוף בכל סט הדגמות. תוצאות מלאות נמצאו עבור שלושת תנאי ההדגמות:
\begin{center}\small
\begin{tabular}{lrrr}
\toprule
מבחן & 0 הדגמות & 4 הדגמות & 12 הדגמות\\\midrule
קשר ישיר & \en{23.5 / 71.5} & \en{87.0 / 90.0} & \en{88.0 / 89.5}\\
שני צעדים & \en{18.5 / 63.5} & \en{62.5 / 78.5} & \en{64.5 / 79.0}\\
ניתוב מחדש & \en{11.5 / 51.0} & \en{39.5 / 49.5} & \en{41.5 / 49.5}\\
עמידות & \en{10.8 / 55.3} & \en{46.3 / 68.2} & \en{47.2 / 68.0}\\\bottomrule
\end{tabular}
\end{center}
בעולם יחסי המשפחה לבדו, ב־12 הדגמות, דיוק \en{candidate} היה $91.9\%$ בקשר הישיר, $86.0\%$ בשני צעדים ו־$46.3\%$ בניתוב מחדש. יש כאן אות התנהגותי מבטיח יותר במבחן הבסיס, אך הוא אינו עובר אוטומטית את הבקרה החזקה. בניית ההדגמות מחדש יכולה לשנות גם את הגרלת הישויות והדאטה; אין לפרש את ההבדל בין גרסאות כהשפעה מבודדת של תיקון יחיד.

\subsection{מגבלות שאותרו בזמן כתיבת התיעוד}
המספרים בגרסאות המאוחרות חושבו מחדש מן השדות \file{free_correct} ו־\file{correct}, ונבדקו $1{,}200$ מזהים ייחודיים בכל ריצה. זו אינה אותה ביקורת מלאה שבוצעה לגרסת ההשלמה הראשונה: לא שוחזר כאן חישוב המודל או אומת כל ציון טוקן מחדש.

נמצאו גם מגבלות קונקרטיות שצריך לתקן לפני שימוש בניתוח הזוגות:
\begin{itemize}
\item ב־\en{v3/v3.1}, $800$ קישורי \file{paired_base_id} בכל ריצה מצביעים למזהה בסיס ללא סיומת סדר, ולכן אינם תואמים למזהי הרשומות. קובצי \file{paired_metrics.json} שנבדקו ריקים; זו אינה תוצאת הצלחה אפס אלא ניתוח שלא הופק.
\item שדות \file{option_order} ו־\file{option_pair_id} נשארו בשם הישן, אך בגרסאות אלה הם מייצגים סדר עובדות, לא סדר אפשרויות גלויות.
\item בתנאי ארבע ושתים־עשרה הדגמות משתנה סט ההדגמות בין שני בני זוג סדר העובדות. לכן אי אפשר לייחס את ההבדל ביניהם לסדר העובדות בלבד.
\item גרסת החלפת השמות אינה שומרת בהכרח מיפוי מפורש כמו הגרסה הקודמת; נדרש שחזור סמנטי לפני השוואת שמות בין גרסאות.
\item המחולל המאוחר משתמש ב־\file{hash()} של פייתון למיקום מסיחים. ללא קיבוע \file{PYTHONHASHSEED}, אותו \en{seed} המוצהר בקוד אינו מבטיח לבדו דאטה זהה בכל תהליך. יש לשמור את הקבצים המדויקים ואת חתימותיהם.
\end{itemize}
מגבלות אלה תועדו כאן; לא שינינו במסגרת כתיבת המסמך את הקוד או את התוצאות ההיסטוריות.

\section{מה למדנו מן התהליך}
\begin{enumerate}
\item \textbf{בנצ'מרק הוא חלק מן המחקר.} גם גרף עם תשובת אמת נכונה יכול להכיל רמז שפותר את המשימה בלי המנגנון שהתכוונו לבדוק.
\item \textbf{יכולת בסיסית תלויה באופן המדידה.} מעבר מכן/לא לשם, הוראה מפורשת, מספר הדגמות ויחס בעל שם שינו את התוצאות. לא כל שינוי בודד בודד ניסויית.
\item \textbf{מודל גדול יותר אינו מבטיח דיוק גבוה יותר בפורמט נתון.} השוואת \en{1B} ל־\en{6.9B} נעשתה לצד שינויים בדאטה ובהקשר, ולכן אינה ניסוי נקי על גודל.
\item \textbf{שמות מומצאים הם פשרה.} הם מצמצמים שליפה של עובדות ספציפיות מהזיכרון, אך עלולים להכביד על טוקניזציה ושליפה מההקשר. ידע קודם על היחס עדיין נחוץ או מועיל.
\item \textbf{דיוק בסיסי אינו עדות מספקת למנגנון.} ניתוב מחדש, עמידות והצלחה מזווגת חשובים יותר מטבלת דיוק יחידה לצורך כניסה למחקר מעגלים.
\item \textbf{תצפית, התערבות ופרשנות הן שלוש רמות שונות.} \en{RI} אינו סיבתיות; \en{patching} הוא התערבות אך אינו לבדו מפת מעגל; מעגל שנמצא צריך לעבור בקרות ותיקוף.
\end{enumerate}

\section{מצב הפרויקט והצעדים שעדיין נדרשים}
\textbf{בוצע:} קריאת ליבה, סקירת כיוונים נוספים, הצעת מתודולוגיה, בחירת מודלים, תשתית הרצה ושחזור, יצירת כמה גרסאות דאטה, פיילוטים על \en{Pythia-1B} ו־\en{Pythia-6.9B}, והפקת ניתוחי טעויות ובקרות.

\textbf{טרם בוצע לפי התיעוד הזמין:} זיהוי \en{SIH} אצלנו באמצעות \en{RI}, ניסויי \en{activation/path patching}, אימות מעגל דליל, מדידת \en{faithfulness}, ניסוי נקודות אימון או אימון משקולות חדש. לכן אין עדיין לכתוב ``גילינו \en{SIC}''.

ההמשך המומלץ הוא לבחור משימת עוגן ברורה, לתקן ולאמת את קישורי הזוגות וההדגמות, להריץ בקרות צינור, ולהקפיא סט פיתוח וסט מבחן עם משפחות חדשות. אין לכוונן שוב ושוב על אותן משפחות ולדווח עליהן בסוף כאילו היו מבחן שלא נראה. לאחר מעבר שער יכולת שנקבע מראש, אפשר להתחיל בסינון רכיבים ובהתערבויות ממוקדות.

אם לא תימצא יכולת עמידה, תוצאה לגיטימית היא תיעוד מגבלות המשימה והמודל או צמצום שאלת המחקר. אין חובה להוכיח מראש שהמנגנון הוא זוג \en{SIH}; ייתכן הסבר אחר או היעדר התנהגות יציבה בתנאים שנבדקו.

\appendix
\section{מפת קבצים לצורך כתיבת העבודה}
כל הנתיבים להלן יחסיים לתיקיית \file{final proj/project}, אלא אם צוין אחרת. הם עוגני שחזור ולא תחליף לשמירת הקבצים הגולמיים.
\begin{itemize}
\item הצעה: \file{../NLP Projects_Proposal.pdf}; ספרות: \file{../papers/}.
\item תכנון מוקדם: \file{data_pipeline_plan.tex}. זהו מסמך תכנון שיש לקרוא בביקורתיות, ולא תיעוד של ניסויים שבוצעו.
\item פיילוט קטן: \file{pythia_pilot/}; תוצאות: \file{results/full_4shot_review/storage/runs/full_4shot/}.
\item כן/לא במודל הגדול: \file{results/full_sdpa_review/}, \file{results/full_sdpa_analysis/}.
\item תיקון הדגמות כן/לא: \file{results/prompt_v3_review/}, \file{results/prompt_v3_analysis/}.
\item השלמה ראשונה: \file{pilot_v2/build_completion_data.py}, \file{results/completion_review/}, \file{results/completion_analysis/}.
\item עומק מפורש: \file{pilot_v2/build_completion_v2_data.py}, \file{results/completion_v2_4shot_v1/}, \file{results/completion_v2_12shot_v1/}.
\item שפה טבעית: \file{pilot_v3/build_data_v3.py}; גרסת קוד העלאה מוקדמת: \file{pilot_v3/upload/build_data_v3.py}.
\item תוצאות שפה טבעית: \file{results/completion_v3_4shot_v1/}, \file{results/completion_v3_12shot_v1/}.
\item תוצאות ההדגמות המאוזנות: \file{results/completion_v3_1_0shot/}, \file{results/completion_v3_1_4shot/}, \file{results/completion_v3_1_12shot/}.
\end{itemize}

\section{כללי דיווח לקראת המאמר הסופי}
לכל טבלה יש לציין מודל וגרסה, דאטה וחתימה, מספר משפחות וקלטים, מספר הדגמות, משמעות המועמדים, מדד וקו בסיס. יש להפריד מדדי אמת מוחלטים ממדדי הצלחה בזוגות, ולציין אילו רכיבים השתנו יחד בין גרסאות. טענה סיבתית תוגבל לניסוי שבוצע. תיעוד של כישלון או של קונפאונד הוא חלק מן התהליך המחקרי, אך אינו תחליף לממצא מכניסטי.

\begin{thebibliography}{9}
\begin{english}
\bibitem{ren} Jie Ren et al. (2024). \emph{Identifying Semantic Induction Heads to Understand In-Context Learning}. Findings of ACL, pp. 6916--6932. \url{https://aclanthology.org/2024.findings-acl.412/}
\bibitem{olsson} Catherine Olsson et al. (2022). \emph{In-context Learning and Induction Heads}. Transformer Circuits Thread. \url{https://transformer-circuits.pub/2022/in-context-learning-and-induction-heads/index.html}
\bibitem{framework} Nelson Elhage et al. (2021). \emph{A Mathematical Framework for Transformer Circuits}. Transformer Circuits Thread. \url{https://transformer-circuits.pub/2021/framework/index.html}
\bibitem{ioi} Kevin Wang et al. (2022). \emph{Interpretability in the Wild: a Circuit for Indirect Object Identification in GPT-2 Small}. \url{https://arxiv.org/abs/2211.00593}
\bibitem{patchguide} Stefan Heimersheim and Neel Nanda (2024). \emph{How to Use and Interpret Activation Patching}. \url{https://arxiv.org/abs/2404.15255}
\end{english}
\end{thebibliography}
\end{document}
